\chapter[Реализация коррекции частотного смещения по пилотам]{Реализация коррекции частотного смещения по пилотам}

\begin{leftbar}
    \href{https://github.com/someotb/sdr}{Ссылка на GitHub}
\end{leftbar}

\section*{Цель практики:}
Реализовать алгоритм частотной коррекции по пилотам, экавализация.

\section*{Краткие теоретические сведения}

\textbf{Эквализация по пилотам}\newline

После удаления циклического префикса и выполнения прямого DFT над
принятым OFDM-символом каждая поднесущая описывается моделью

\begin{equation}
    Y[k] = H[k]\,X[k] + W[k], \quad k = 0,\dots,N-1,
    \label{eq:channel_model}
\end{equation}

где $H[k]$ — комплексный коэффициент передачи канала на $k$-й поднесущей
(амплитудное затухание и фазовый сдвиг, вызванные многолучевым
распространением), $X[k]$ — переданный символ, $W[k]$ — аддитивный шум.
Поскольку канал частотно-селективен, $H[k]$ различен для разных
поднесущих, и для восстановления $X[k]$ его необходимо оценить и
скомпенсировать отдельно на каждой поднесущей.

\vspace{20pt}

\textbf{Оценка канала по пилотам}\newline

На части поднесущих (пилотных, с индексами $k \in \mathcal{P}$)
передаются символы $X_p[k]$, заранее известные приёмнику. Это позволяет
получить прямую оценку канала на этих поднесущих:

\begin{equation}
    \hat{H}[k] = \frac{Y[k]}{X_p[k]}, \quad k \in \mathcal{P}.
\end{equation}

Поскольку пилоты занимают не все поднесущие (иначе не осталось бы места
под полезные данные), оценка канала для промежуточных поднесущих
$k \notin \mathcal{P}$ получается интерполяцией между соседними
пилотными оценками (линейной или иного порядка).

\vspace{20pt}

\textbf{Компенсация искажений (эквализация)}\newline

Зная оценку $\hat{H}[k]$ для каждой поднесущей, искажение канала
устраняется делением принятого символа на эту оценку:

\begin{equation}
    \hat{X}[k] = \frac{Y[k]}{\hat{H}[k]}.
\end{equation}

Эта операция одновременно компенсирует и амплитудное затухание, и
фазовый сдвиг, внесённые каналом на данной конкретной поднесущей, что
возможно только благодаря диагонализации канала в частотной области,
обеспеченной циклическим префиксом.

\vspace{20pt}

\textbf{Отличие от коррекции CFO}\newline

Эквализация по пилотам и коррекция частотного смещения решают разные
задачи. CFO вносит фазовый сдвиг, одинаково нарастающий во времени для
всех поднесущих одновременно и не зависящий от частотной характеристики
канала, в то время как $H[k]$ отражает избирательные по частоте
искажения, вносимые самим каналом. Если CFO не скомпенсировано до FFT,
вносимая им фазовая ошибка исказит и саму оценку $\hat{H}[k]$ по
пилотам, поскольку эквалайзер интерпретирует эту ошибку как часть
характеристики канала.

\section*{Выполнение}

Была написана функция эквализации:

\begin{lstlisting}[language=C++,caption=Фунция эквализации]
void equalization(std::vector<std::complex<float>> &in_signal, sharedData &sd, std::vector<std::complex<float>> &out_signal)
{
    out_signal.clear();
    out_signal.reserve(in_signal.size());

    int subcarrar = sd.subcarrier;
    int num_symbols = in_signal.size() / subcarrar;

    for (int i = 0; i < num_symbols; ++i)
    {
        std::vector<std::complex<float>> H(subcarrar, {0.0f, 0.0f});
        std::vector<std::complex<float>> equalized(subcarrar, {0.0f, 0.0f});
        const std::complex<float> known_pilot_conj(1.0f / std::sqrt(2.0f), 0.0f / std::sqrt(2.0f));

        for (int p_idx : sd.pilot_idxs)
        {
            std::complex<float> received_pilot = in_signal[i * subcarrar + p_idx];
            H[p_idx] = received_pilot * known_pilot_conj;
        }

        for (size_t p = 0; p + 1 < sd.pilot_idxs.size(); ++p)
        {
            int k1 = sd.pilot_idxs[p];
            int k2 = sd.pilot_idxs[p + 1];

            std::complex<float> H1 = H[k1];
            std::complex<float> H2 = H[k2];

            float phi1 = std::arg(H1);
            float phi2 = std::arg(H2);

            float delta_phi = phi2 - phi1;

            if (delta_phi > M_PIf)
                delta_phi -= 2 * M_PI;

            if (delta_phi < -M_PIf)
                delta_phi += 2 * M_PI;

            for (int k = k1 + 1; k < k2; ++k)
            {
                if (sd.is_zeros[k])
                    continue;

                float alpha = float(k - k1) / float(k2 - k1);
                H[k] = std::polar(std::abs(H1) + alpha * (std::abs(H2) - std::abs(H1)), phi1 + alpha * delta_phi);
            }
        }

        for (int k = 0; k < sd.pilot_idxs.front(); ++k)
            H[k] = H[sd.pilot_idxs.front()];

        for (int k = sd.pilot_idxs.back() + 1; k < subcarrar; ++k)
            H[k] = H[sd.pilot_idxs.back()];

        {
            std::lock_guard<std::mutex> lock(sd.sync.channel_estimation_mutex);
            if (sd.channel_estimation.size() >= (size_t)sd.mtu)
                sd.channel_estimation.erase(sd.channel_estimation.begin(),
                                                  sd.channel_estimation.begin() + sd.subcarrier);
            for (const auto& h : H)
                sd.channel_estimation.push_back(h);
        }

        for (int k = 0; k < subcarrar; ++k)
        {
            if (sd.is_zeros[k] || std::abs(H[k]) < 1e-6f)
                equalized[k] = {0, 0};
            else
                equalized[k] = in_signal[i * subcarrar + k] / H[k];
        }

        for (size_t k = 0; k < equalized.size(); ++k)
        {
            if (sd.is_zeros[k] || sd.is_pilot[k])
                continue;

            out_signal.push_back(equalized[k]);
        }
    }
}
\end{lstlisting}

\textbf{Функция эквализации}\newline

Функция \texttt{equalization} выполняет коррекцию частотно-избирательных
искажений канала для каждого OFDM-символа входного буфера отдельно.

Для каждого символа сначала на пилотных поднесущих вычисляется прямая
оценка канала $H[k] = Y[k] \cdot \overline{X_p}[k]$, где $Y[k]$ —
принятое значение пилота, $\overline{X_p}[k]$ — сопряжённое известное
значение пилота. Далее оценка канала для поднесущих между соседними
пилотами получается линейной интерполяцией отдельно по модулю и по
фазе $H[k]$, с корректной обработкой перехода фазы через $\pm\pi$.
Поднесущие до первого и после последнего пилота получают оценку,
равную ближайшему крайнему пилоту (экстраполяция константой).

Полученная оценка канала сохраняется в общий буфер \texttt{sd.channel\_estimation}
(с ограничением по размеру \texttt{mtu}) для последующего использования,
например при построении графиков АЧХ/ФЧХ канала.

Компенсация искажений выполняется делением принятого символа на оценку
канала: $\hat{X}[k] = Y[k] / H[k]$. Нулевые (защитные) поднесущие и
поднесущие с вырожденной оценкой канала ($|H[k]| < 10^{-6}$)
принудительно обнуляются во избежание деления на близкое к нулю число.

На выходе из каждого символа исключаются пилотные и нулевые поднесущие —
в выходной буфер попадают только эквализованные символы данных.

Вот как выглядит работа CFO и эквализации между 2-мя Adalm Pluto SDR:

\begin{figure}[H]
    \centering
    \includegraphics[width=0.6\textwidth]{nothing.jpg}
    \caption{Передача без CFO и EQ}
\end{figure}

\begin{figure}[H]
    \centering
    \includegraphics[width=0.6\textwidth]{add_cfo.jpg}
    \caption{Передача c CFO, но без EQ}
\end{figure}

\begin{figure}[H]
    \centering
    \includegraphics[width=0.6\textwidth]{add_eq.jpg}
    \caption{Передача с CFO и EQ}
\end{figure}

\section*{Результаты практики}

Была написана программа, которая реализует полноценную передачу OFDM символов, с CFO и EQ. Поддерживаемые модуляции:
\begin{itemize}
	\item BPSK
	\item QPSK
	\item QAM16
	\item QAM64
	\item QAM256
\end{itemize}

Также была создана панель с отладочной информация, которая показывает:
\begin{itemize}
	\item Задержку dsp уровня, которая показывает сколько милисекунд занимает обрабтка сэмплов
	\item График с пиком корреляции Zadoff-Chu
	\item График нахождения пиков Zadoff-Chu
	\item График коррекции CFO
	\item График оценки канала, которая берется из EQ
\end{itemize}

\begin{figure}[H]
    \centering
    \includegraphics[width=0.6\textwidth]{latency.jpg}
    \caption{Задержка DSP уровня}
\end{figure}

\begin{figure}[H]
    \centering
    \includegraphics[width=0.6\textwidth]{channel_est.jpg}
    \caption{Оценка канала}
\end{figure}

\begin{figure}[H]
    \centering
    \includegraphics[width=0.6\textwidth]{cfo_corr.jpg}
    \caption{График CFO}
\end{figure}

Пример передачи QAM16 между 2-х Adalm Pluto SDR

\begin{figure}[H]
    \centering
    \includegraphics[width=1.0\textwidth]{ready.jpg}
    \caption{Передача QAM16(2 SDR)}
\end{figure}

\section*{Вывод}

В ходе практики была реализована полная цепочка передачи и приёма
OFDM-сигнала: формирование символов из потока данных с поддержкой
модуляций BPSK, QPSK, QAM16, QAM64 и QAM256, детектирование границ
кадра по корреляции с последовательностью Задова-Чу, оценка и
коррекция частотного смещения по циклическому префиксу, оценка канала
и эквализация по пилотным поднесущим.

Работоспособность реализованных алгоритмов подтверждена не только на
модельных данных, но и на реальном радиоканале между двумя SDR
Adalm Pluto — сигнал с модуляцией QAM16 был успешно передан, принят и
демодулирован без ошибок, что отражено на приведённых графиках оценки
канала и коррекции CFO.

Отдельно была разработана отладочная панель, отображающая ключевые
метрики DSP-цепочки в реальном времени: задержку обработки на уровне
DSP, пик корреляции и положение синхронизирующей последовательности
Задова-Чу, динамику коррекции CFO и текущую оценку канала. Эти
инструменты подтвердили, что каждый этап обработки — синхронизация,
коррекция частоты, эквализация — работает согласованно и вносит
измеримый, ожидаемый вклад в качество приёма.

Практика позволила перейти от теоретического понимания OFDM к
работающей системе, где каждый этап обработки сигнала — от
формирования символа до восстановления данных — реализован,
протестирован и проверен на реальном оборудовании.
